\section*{Introduction}
\label{sec:intro}

\addcontentsline{toc}{section}{Introduction}

The purpose of the present paper is to introduce the Mukai pairing on
the Hochschild homology of smooth, proper spaces.  This pairing is the
natural analogue, in the context of Hochschild theory, of the
Poincar\'e pairing on the singular cohomology of smooth manifolds.

Our approach is categorical.  We start with a geometric category,
whose objects will be called {\em spaces}.  For a space $X$ we define
its Hochschild homology which is a graded vector space $\HH_\blob(X)$
equipped with the non-degenerate \emph{Mukai pairing}.  We show that
this structure satisfies a number of properties, the most important of
which are {\em functoriality} and {\em adjointness}.

The advantage of the categorical approach is that the techniques we
develop apply in a wide variety of geometric situations, as long as an
analogue of Serre duality is satisfied.  Examples of categories for
which our results apply include compact complex manifolds, proper
smooth algebraic varieties, proper Deligne-Mumford stacks for which
Serre duality holds, representations of a fixed finite group, and
compact ``twisted spaces'' in the sense of~\cite{Cal}.  We expect the
same construction to work for categories of Landau-Ginzburg
models~\cite{Orl}, but at the moment it is not known whether Serre
duality holds for these.

\subsection*{The Hochschild structure}
\label{subsec:hochstr}
In order to define the Hochschild structure of a space we need
notation for certain special kernels which will play a fundamental
role in what follows.  For a space $X$, denote by $\Id_X$ and
$\SerK_X^{-1}$ the objects of $\D(X\times X)$ given by
\[ \Id_X := \Delta_* \cO_X \quad\mbox{and}\quad \SerK_X^{-1} :=
\Delta_* \omega_X^{-1}[-\dim X], \] 
where $\Delta:X \ra X\times X$ is the diagonal map, and
$\omega_X^{-1}$ is the anti-canonical line bundle of $X$.  When
regarded as kernels, these objects induce the identity functor and the
inverse of the Serre functor on $\D(X)$, respectively.  We shall see
in the sequel that $\Id_X$ can be regarded as the identity 1-morphism
of $X$ in a certain 2-category $\gVar$.

The {\em Hochschild structure} of the space $X$ then consists of the
following data:
\begin{itemize}
\item the graded ring $\HH^\blob(X)$, the Hochschild cohomology ring of
  $X$, whose $i$-th graded piece is defined as
  \[ \HH^i(X) := \Hom_{\D(X\times X)}^i(\Id_X, \Id_X); \]
\item the graded left $\HH^\blob(X)$-module $\HH_\blob(X)$, the
  Hochschild homology module of $X$, defined as
  \[ \HH_i(X) := \Hom_{\D(X\times X)}^{-i}(\SerK_X^{-1}, \Id_X); \]
\item a non-degenerate graded pairing $\MP{{-}}{{-}}$ on
  $\HH_\blob(X)$, the {\em generalized Mukai pairing}.
\end{itemize}

The above definitions of Hochschild homology and cohomology agree with
the usual ones for quasi-projective schemes (see~\cite{CalHH2}).  The
pairing is named after Mukai, who was the first to introduce a pairing
satisfying the main properties below, on the total cohomology of
complex K3 surfaces~\cite{MukK3}.

\subsection*{Properties of the Mukai pairing}
The actual definition of the Mukai pairing is quite complicated and is
given in Section~\ref{sec:adjoint}.  We can, however, extricate the
fundamental properties of Hochschild homology and of the Mukai
pairing.
\begin{Description}
\item[Functoriality] Integral kernels induce, in a functorial way,
  linear maps on Hochschild homology. Explicitly, to any integral
  kernel \( \Phi\in \D(X\times Y) \) we associate, in
  Section~\ref{def:defmaps}, a linear map of graded vector spaces
  \[ \Phi_* \colon \HH_\blob(X) \ra \HH_\blob(Y), \]
  and this association is functorial with respect to composition of
  integral kernels (Theorem~\ref{thm:functorial}).
\item[{Adjointness}] For any adjoint pair of integral kernels
  $\Psi\adjoint \Phi$, the induced maps on homology are themselves
  adjoint with respect to the Mukai pairing:
  \[ \MP{\Psi_* v}{w} = \MP{v}{\Phi_* w}\]
  for $v\in \HH_\blob(Y)$, $w \in \HH_\blob(X)$
  (Theorem~\ref{thm:adjointness}).
\end{Description}
The following are then consequences of the above basic properties:
\begin{Description}
\item[{Chern character}] In all geometric situations there is a
  naturally defined object $\bone \in \HH_0(\pt)$.  An element $\cE$
  in $\D(X)$ can be thought of as the kernel of an integral transform
  $\pt \ra X$, and using functoriality of homology we define a Chern
  character map
  \[ \ch\colon\K_0(X) \ra \HH_0(X), \quad \ch(\cE) = \cE_*(\bone).  \]
  For a smooth proper variety the Hochschild-Kostant-Rosenberg
  isomorphism identifies $\HH_0(X)$ and $\bigoplus_p \H^{p,p}(X)$; our
  definition of the Chern character matches the usual one under this
  identification~\cite{CalHH2}.
\item[Semi-Hirzebruch-Riemann-Roch Theorem] For $\cE, \cF \in \D(X)$
  we have
  \[ \MP{\ch(\cE)}{\ch(\cF)} = \chi(\cE, \cF) = \sum_i (-1)^i \dim
  \Ext_X^i(\cE, \cF). \]
\item[Cardy Condition] The Hochschild structure appears naturally in
  the context of open-closed topological quantum field theories
  (TQFTs).  The Riemann-Roch theorem above is a particular case of a
  standard constraint in these theories, the Cardy Condition.  We
  briefly discuss open-closed TQFTs, and we argue that the natural
  statement of the Cardy Condition in the B-model open-closed TQFT is
  always satisfied, even for spaces which are not Calabi-Yau
  (Theorem~\ref{thm:cardy}).
\end{Description}

\subsection*{The 2-categorical perspective}
In order to describe the functoriality of Hochschild homology it is
useful to take a macroscopic point of view using a 2-category called
$\gVar$.  One way to think of this 2-category is as something half-way
between the usual category consisting of spaces and maps, and $\gCat$,
the 2-category of (derived) categories, functors and natural
transformations.  The 2-category $\gVar$ has spaces as its objects,
has objects of the derived category $\D(X\times Y)$ --- considered as
integral kernels --- as its 1-morphisms from $X$ to $Y$, and has
morphisms in the derived category as its
2-morphisms.

One consequence of thinking of spaces in this 2-category is that
whereas in the usual category of spaces and maps two spaces are
equivalent if they are isomorphic, in $\gVar$ two spaces are
equivalent precisely when they are Fourier-Mukai partners.  This is
the correct notion of equivalence in many circumstances, thus making
$\gVar$ an appropriate context in which to work.

This point of view is analogous to the situation in Morita theory in
which the appropriate place to work is not the category of algebras
and algebra morphisms, but rather the 2-category of algebras,
bimodules and bimodule morphisms.  In this 2-category two algebras are
equivalent precisely when they are Morita equivalent, which again is
the pertinent notion of equivalence in many situations.

Many facts about integral transforms can be stated very elegantly as
facts about the 2-category $\gVar$.  For example, the fact that every
integral transform between derived categories has both a left and
right adjoint is an immediate consequence of the more precise fact ---
proved exactly the same way --- that every integral kernel has both a
left and right adjoint in $\gVar$.  Here the definition of an adjoint
pair of 1-morphisms in a 2-category is obtained from one of the
standard definitions of an adjoint pair of functors by everywhere replacing the
word `functor' by the word `1-morphism' and the words `natural
transformation' by the word `2-morphism'.

The Hochschild cohomology of a space $X$ has a very natural
description in terms of the 2-category $\gVar$: it is the ``second
homotopy group of $\gVar$ based at $X$'', which means that it is
$\THom_\gVar(\Id_X,\Id_X)$, the set of 2-morphisms from the identity
1-morphism at $X$ to itself.  Unpacking this definition for $\gVar$
one obtains precisely $\Ext^\blob_{X\times X}(\cO_\Delta,\cO_\Delta)$,
one of the standard definitions of Hochschild cohomology.  By analogy
with homotopy groups, given a kernel $\Phi\colon X\to Y$, i.e., a
``path'' in $\gVar$, one might expect an induced map $\HH^\blob(X)\to
\HH^\blob(Y)$ obtained by ``conjugating with $\Phi$''.  However, this
does not work, as the analogue of the ``inverse path to $\Phi$''
needed is a simultaneous left and right adjoint of $\Phi$, and such a
thing does not exist in general as the left and right adjoints of
$\Phi$ differ by the \emph{Serre kernels} of $X$ and $Y$.

The Hochschild homology $\HH_\blob(X)$ of a space $X$ can be given a
similar natural definition in terms of $\gVar$ --- it is
$\THom_\gVar(\SerK^{-1}_X,\Id_X)$ the set of 2-morphisms from the
inverse Serre kernel of $X$ to the identity 1-morphism at $X$.
In this case, the idea of ``conjugating by a kernel $\Phi\colon X \to
Y$'' does work as the Serre kernel in the definition exactly
compensates the discrepancy between the left and right adjoints of
$\Phi$.

The functoriality of Hochschild homology can be expressed by saying
that $\HH_\blob$ is a functor into the category of vector spaces from
the Grothen\-dieck category of the 2-category $\gVar$ (i.e., the
analogue of the Grothendieck group of a 1-category) whose objects are
spaces and whose morphisms are \emph{isomorphism classes} of kernels.
One aspect of this which we do not examine here is related to the fact
that this Grothendieck category is actually a monoidal category with
certain kinds of duals for objects and morphisms, and that Hochschild
homology is a monoidal functor.  The Mukai pairing is then a
manifestation of the fact that spaces are self-dual in this
Grothendieck category.  Details will have to appear elsewhere.

There is an alternative categorical approach to defining Hoch\-schild
homology and cohomology.  This approach uses the notion of enhanced
triangulated categories of Bondal and Kapranov~\cite{BonKap}, which
are triangulated categories arising as homotopy categories of
differential-graded (dg) categories.  In~\cite{Toe}, To\"en argued that
the Hochschild cohomology $\HH^\blob(X)$ of a space $X$ can be
regarded as the cohomology of the dg-algebra of dg-natural
transformations of the identity functor on the dg-enhancement of
$\D(X)$.  It seems reasonable to expect that a similar construction
can be used to define the Hochschild homology $\HH_\blob(X)$ as
dg-natural transformations from the inverse of the Serre functor to
the identity.  However, since the theory of Serre functors for
dg-categories is not yet fully developed, we chose to use the language
of the 2-category $\gVar$, where all our results can be made precise.

\subsection*{String diagram notation}
As 2-categories are fundamental to the functoriality, and they are
fundamentally 2-dimensional creatures, we adopt a 2-dimensional
notation.  The most apt notation in this situation appears to be that
of \emph{string diagrams}, which generalizes the standard notation
used for monoidal categories in quantum topology.  String diagrams are
Poincar\'e dual to the usual arrow diagrams for 2-categories.  The
reader unfamiliar with these ideas should be aware that the pictures
scattered through this paper form rigorous notation and are not just
mnemonics.

\subsection*{Note} This paper supersedes the unpublished paper
\cite{CalMPI}, in which it was stated that hopefully the correct
categorical context could be found for the results therein.  This
paper is supposed to provide the appropriate context.

\subsection*{Synopsis} The paper is structured as follows.  The first
section is devoted to the study of integral transforms and of the
2-category $\gVar$.  In the next section we review the Serre functor and
Serre trace on the derived category, and we use these in
Section~\ref{sec:adjoints} to study adjoint kernels in $\gVar$.  In
Section~\ref{sec:defmaps} we introduce the maps between Hochschild
homology groups associated to a kernel, and we examine their
functorial properties.  The Mukai pairing is defined in
Section~\ref{sec:adjoint}, where we also prove its compatibility with
adjoint functors.  In Section~\ref{sec:chernchar} we define the Chern
character and we prove the Semi-Hirzebruch-Riemann-Roch theorem.  We conclude with
Section~\ref{sec:cardy} where we review open-closed TQFTs, and we
discuss the Cardy Condition.  An appendix contains some of the more
technical proofs.

\subsection*{Notation} Throughout this paper $\k$ will denote an
algebraically closed field of characteristic zero and $\D(X)$ will
denote the bounded derived category of coherent sheaves on $X$.
Categories will be denoted by bold letters, such as $\C$, and the
names of 2-categories will have a script initial letter, such as
$\gVar$.


 \subsection*{The base category of spaces}
We fix for the remainder of the paper a geometric category, whose
objects we shall call \emph{spaces}.  It is beyond the purpose of this
paper to list the axioms that this category needs to satisfy, but the
following categories can be used:
\begin{itemize}
\item smooth projective schemes over $\k$;
\item smooth proper Deligne-Mumford stacks over $\k$;
\item smooth projective schemes over $\k$, with an action of a fixed
  finite group $G$, along with $G$-equivariant morphisms;
\item twisted spaces in the sense of~\cite{Cal}, i.e., smooth
  projective schemes over $\k$, enriched with a sheaf of Azumaya
  algebras.
\end{itemize}
For any space $X$ as above, the category of coherent sheaves on $X$
makes sense, and the standard functors (push-forward, pull-back,
sheaf-hom, {etc}.)  are defined and satisfy the usual compatibility
relations.

\subsection*{Acknowledgments} We have greatly benefited from
conversations with Jonathan Block, Tom Bridgeland, Andrew Kresch,
Aaron Lauda, Eyal Markman, Mircea Musta\c t\u a, Tony Pantev, Justin
Roberts, Justin Sawon, and Sarah Witherspoon. Many of the ideas in
this work were inspired by an effort to decipher the little known but
excellent work~\cite{Mar} of Nikita Markarian. Greg Moore and Tom
Bridgeland suggested the connection between the Semi-Hirzebruch-Riemann-Roch
theorem and the Cardy Condition in physics.

AC's initial work on this project was supported by an NSF
post-doctoral fellowship, and by travel grants and hospitality from
the University of Pennsylvania, the University of Salamanca, Spain,
and the Newton Institute in Cambridge, England.  AC's current work is
supported by the National Science Foundation under Grant No.\
DMS-0556042.  SW has been supported by a WUN travel bursary and a
Royal Society Conference grant.
